\clearpage
\fancyhead[L]{\small A NEW AID.}
\SourceLeaf{358}
\pdfbookmark[0]{Göttingische gelehrte Anzeigen. October 30, 1837.}{anzeige-1837-10-30}
\begin{center}
Göttingische gelehrte Anzeigen. October 30, 1837.
\end{center}

The lecture given by Court Councillor \textsc{Gauss} at the meeting of the Royal Society on September 19 concerns
\begin{center}
\textit{a new aid for magnetic observations.}
\end{center}

Magnetic declination is one of the elements in which terrestrial magnetism manifests itself. It was the first to occupy observers, and even since naturalists have turned their attention to the other elements, declination has continued to take precedence over them in several respects. One of the most interesting investigations in terrestrial magnetism concerns the remarkable irregular disturbances of this force, which nevertheless act simultaneously and uniformly across an entire continent. A separate association of observers has come together in recent years to study them, and the other two elements have so far been entirely excluded from that investigation. The reason for preferring declination to the other elements lies less in the considerable practical importance of knowing the declination for navigators, geodesists and mine surveyors than in the state of the available instruments. For declination they leave nothing to be desired, while for the other elements they are still far behind. The magnetometer introduced some years ago does indeed serve to determine both declination and horizontal intensity. In fact, the problem of deter\SourceLeaf{359}mining the latter in absolute measure was what first prompted its construction. Yet it still does not solve the problem in \textit{every} respect. By its nature it can give a precise value only for the mean intensity over a certain period, and although this period. is to some extent ours to choose, it must not be made too short, because that would change the precision and reliability of the result as well. The magnetometer therefore cannot be used at all to follow changes in intensity over short intervals. The new apparatus discussed in the lecture is intended to fill this gap for intensity and make observations of this element as easy and reliable as the magnetometer makes observations of declination.

The limited space in these pages does not allow a description of the apparatus that Court Councillor \textsc{Gauss} commissioned. Such a description is all the less necessary because the lecture itself will soon appear in print. The instrument has been suspended in the astronomical observatory here for several months, and changes in intensity have already been observed on it during two magnetic observing sessions, for 24 hours each time. If the changes in declination observed at the same time in the magnetic observatory are plotted with them in a single diagram, the wonderful interplay of the magnetic disturbances stands out vividly in a new and distinctive way. Once similar simultaneous observations are made at several widely separated places, we can confidently expect soon to learn more fully where the causes of these puzzling phenomena are located. Intensity shows irregular disturbances as frequent and as considerable as those of declination. Both also show regular changes connected with the time of day. In intensity, however, these changes take a somewhat different course, as far as can be judged from a month of daily records at fixed hours. In our region the westerly declination increases from 7 or 8 in the morning until one or two hours after noon, and then recedes. Intensity, meanwhile, decreases in the early morning hours, but reaches its minimum as early as one or two hours before noon, just when declination is increas\SourceLeaf{360}ing most rapidly. It scarcely needs pointing out that on individual days irregular disturbances often obscure this regular pattern completely, and that more precise determinations can come only from observations continued over a long period.

The design of the apparatus allows it to be used for many quite different purposes besides its main one. It solves a problem that people had repeatedly tried to solve without success, namely how to display the hourly and irregular changes in declination on an enlarged scale. As the apparatus is arranged at present, the enlargement is tenfold. A change in declination that shows as 50 scale divisions on the magnetometer in the magnetic observatory appears here as 500 divisions. At the last magnetic observing session (September 30), 8 hours of simultaneous observations with both instruments confirmed this most satisfactorily.

As is well known, combining the magnetometer with a multiplier coil makes a very sensitive galvanometer. The new apparatus can be used in the same way, but its sensitivity exceeds that of the magnetometer in exactly the same ratio we gave for changes in declination. The new apparatus therefore permits precise measurement even of the weakest galvanic currents. It is often astonishing to see these currents set the 25-pound magnetic bar inside it in such considerable motion. With currents excited thermogalvanically, this very clearly refutes the mistaken view of many physicists that such currents cannot pass through a circuit of considerable length. However long the circuit, it does not extinguish these currents. It only weakens them, in exactly the same ratio as with other means of excitation. With a thermogalvanic apparatus of a particular design, merely touching the junctions with a finger produces a galvanic current. Even if it has to pass through a circuit almost two Meilen long, made mostly of very thin wire, it can still be detected by very large deflections of the magnetic bar.

When ordinary frictional electricity is passed through the multiplier coil, its electromagnetic effects are among the weakest, hard to detect and even harder to measure. As is well known, the
\SourceLeaf{361}existence of such effects was first discovered by \textsc{Colladon} and later. confirmed by \textsc{Faraday}. Instead of discharging a powerful electric battery through the conducting wire, as those physicists did, Court Councillor \textsc{Gauss} observed the effect of frictional electricity while an electrical machine set up in the physics cabinet was kept turning. Its conductor and rubbing pad were connected to the ends of the large circuit running to the astronomical observatory. This circuit included the multiplier coil surrounding the magnetic bar of the new apparatus. The bar was thereby held at a deflection of 144 scale divisions or 51 minutes of arc, positive or negative according to the direction in which the electricity passed through the coil. The length of wire in the circuit was about 13000 feet. What is particularly remarkable is that lengthening the circuit to nearly a whole Meile by adding other wire brought about no reduction at all in the electromagnetic effect. In this respect the current from the electrical machine behaves differently from galvanic currents excited hydrogalvanically, thermogalvanically or by induction. Their intensity, as measured by the magnetic effect, decreases as the length of the circuit that closes the current path increases. Far from proving an essential difference in the nature of these currents, that phenomenon confirms their sameness and the theory according to which a difference in the intensity of two currents amounts simply to a difference in the quantity of electricity passing through any cross section of the conductor in a given time. For given electromotive forces of the kinds last mentioned, less electricity is set in motion as the resistance offered by the longer circuit increases. But in the experiment described above, \textit{all} the electricity that jumped from the machine to the conductor in sparks had to pass through the whole circuit to balance the opposite electricity of the rubbing pad, whether the circuit was short or long (provided it was sufficiently insulated). The quantity of electricity passing through any cross section of the conducting wire in a given time therefore depended entirely on the operation of the machine, regardless of the length of the circuit.

In most of the experiments mentioned here, the galvanic or electric current had to pass through the large circuit laid in 1833 between the astronomical observatory and the physics \SourceLeaf{362}cabinet. The wire in the open air alone is almost 7000 feet long. The main purpose of this installation is to conduct large-scale physical investigations into the laws of galvanic currents. Yet from the outset it was also used for numerous experiments in electromagnetic telegraphy, which succeeded most satisfactorily even with whole words and short phrases. This was already mentioned in these pages in the first report on the establishment of the magnetic observatory here (Gött. gel. Anz. 1834, August 9). In 1835 the original method was replaced by another that Court Councillor \textsc{Gauss} reached by considering the laws of induction, and which is far preferable to the first. It is in precisely this mode of telegraphy that the new apparatus under discussion has a significant advantage over the ordinary magnetometer. This prompted Court Councillor \textsc{Gauss} to describe in his lecture the main features of this method, which had yet to be mentioned publicly, and to explain what it achieves. We must pass over this here because of the limited space. For the same reason we mention only briefly another new device intended to bring any untimely oscillation of a magnetic needle that interferes with a particular operation to rest on its own in a short time. It is therefore called a damper. The device can be used equally well with the new apparatus or with the magnetometer and is of substantial use both in the mode of telegraphy just mentioned and in many other magnetic measurements.

